\documentclass[10pt]{extarticle}

\usepackage{custom}

\title{Project List 2}
\author{33-467: Astrophysics of the Stars and Galaxy}
\date{Fall 2023}

\begin{document}
\maketitle

\definecolor{tcol_CNT1}{HTML}{EDEADE} % First color for Contents
\definecolor{tcol_CNT2}{HTML}{EDEADE} % Second color for Contents
\definecolor{tcol_CNV1}{HTML}{8E44AD} % First color for Conventions
\definecolor{tcol_CNV2}{HTML}{A10B49} % First color for Conventions

\begin{tcolorbox}[enhanced,
    fonttitle=\fontsize{20}{24}\sffamily\bfseries\selectfont, coltitle=black,
    fontupper=\sffamily, interior style={left color=tcol_CNT1!80,right
    color=tcol_CNT2!80}, frame style={left color=tcol_CNT1!60!black,right
    color=tcol_CNT2!60!black}, attach boxed title to top center={yshift=10pt},
    boxed title style={frame hidden, interior style={left color=tcol_CNT1,right
    color=tcol_CNT2}, frame style={left color=tcol_CNT1!60!black,right
    color=tcol_CNT2!60!black}, height=24pt,bean arc,drop fuzzy shadow },
    top=2mm,bottom=2mm,left=2mm,right=2mm, before skip=20mm,after skip=20mm,
    drop fuzzy shadow,breakable]
%
\makeatletter
\@starttoc{toc}
\makeatother
\end{tcolorbox}

\section{Binary evolution with COSMIC}

\subsection{Orbital Evolution of Binaries in Roche-lobe Overflow}
This project has both analytic and computational components. In class, we will
cover how a binary star's orbit responds to mass exchange through Roche-lobe
overflow from one star to its companion. Your first task is to reproduce the
derivation under the assumption that the mass transfer is fully conservative. In
the Universe, though, it is not required that mass transfer remain fully
conservative. In fact, the amount of mass a star can accrete through Roche-lobe
overflow is an open research question! 

Since we don't know the answer for how conservative mass transfer can be, we
parameterize this uncertainty and calculate how different assumptions for how
much mass is accreted affect the future evolution of the binary. You should
introduce a tuning parameter (often called $\beta$) that represents how much
mass is accreted. Before running any code, write down some predictions for how
different values for the accretion fraction will alter the future evolution of a
binary. Once you have your predictions, install COSMIC and explore the
\href{https://cosmic-popsynth.github.io/COSMIC/}{COSMIC documentation} and
\href{https://ui.adsabs.harvard.edu/abs/2020ApJ...898...71B/abstract}{COSMIC
code paper}. Next, initialize a binary that will undergo stable mass transfer
through Roche-lobe overflow (this will require some guess and check). Using
different assumptions for the accretion fraction (`f\_acc'), determine whether
your predictions were correct. 

\subsection{The Wild World of Binary-Star Evolution}
There are \emph{many} different uncertainties associated with binary evolution
due to different interactions between binary components that can occur based on
the initial properties of the system and assumptions made for how each
interaction should be treated. This means that there are \emph{a lot} of
opportunities for exploration. You should explore the
\href{https://cosmic-popsynth.github.io/COSMIC/}{COSMIC documentation} and
\href{https://ui.adsabs.harvard.edu/abs/2020ApJ...898...71B/abstract}{COSMIC
code paper} for a summary of the different assumptions made in the code as well
as the
\href{https://www.astro.ru.nl/~sluys/Binaries/files/BinaryEvolutionNutshell_letter.pdf}{Binary
Evolution in a Nutshell} notes to find an aspect of binary star interactions
that you find most interesting. Maybe it's studying how different assumptions
for wind mass loss affects binary populations? Maybe it's investigating how
compact object populations depend on the strength of supernova natal kicks?
Maybe it's something else!

Whatever topic you choose should not be exactly the same as any other projects,
so you may need to coordinate with your colleagues. Once you have a topic to
study, you should determine which assumptions are set by flags in COSMIC that
affect that topic as well as determine which binary populations will show the
largest changes by using different assumptions. You should then summarize how
you expect each assumption you make to alter the binary population from its
initial configuration at ZAMS. Finally, you should evolve large populations of
binaries with COSMIC using the flags you've chosen and make figures which
illustrate how each assumption impacts the population. 


\section{Finding Dark Matter with Stellar Streams}
One proposed form of dark matter is the $\Lambda$ Cold Dark Matter
($\Lambda$CDM) model which is expected to result in small-scale dark matter
halos orbiting along with stars in galaxies. Stellar streams, made up of tidally
disrupted globular clusters and dwarf galaxies, are \emph{amazing} tracers of
small-scale gravitational overdensities like dark matter clumps.

In this project, you will follow along with a
\href{https://datacarpentry.org/astronomy-python/}{Data Carpentry Tutorial} to
reporduce the results of
\href{https://ui.adsabs.harvard.edu/abs/2018ApJ...863L..20P/abstract}{Price-Whelan
et al. 2019} which shows that stars in the stellar stream Gaia GD-1 may have
been perturbed due to a dark matter clump. This project will involve building
skills in astroquery, joining different data sets, and selecting regions of
interest in a dataset. By the end of the project, you should have figures which
show the Gaia GD-1 stream and illustrate why the authors argue that the stream
has been perturbed by a small-scale dense object that could be a dark matter
subhalo. As part of this project, you should also summarize what stellar streams
are, what the $\Lambda$CDM model predicts and why it is favored by the
scientific community, and how future observations could aid in discoveries of
dark matter clumps.



\section{Black hole kicks with Gaia BH1 and Gaia BH2}
The recent discoveries of
\href{https://ui.adsabs.harvard.edu/abs/2023MNRAS.518.1057E/abstract}{Gaia BH1}
and \href{https://ui.adsabs.harvard.edu/abs/2023MNRAS.521.4323E/abstract}{Gaia
BH2} were some of the most exciting discoveries in 2023 (at least according to
Prof. Breivik). Since there is no accretion in either binary, the last thing to
alter the orbit of the binary and its trajectory in the Galaxy was the natal
kick that may or may not have been imparted during the formation of each black
hole. This means that the current orbit of each binary in the Galaxy places a
pretty strong constraint on the black hole kicks!

For this project, you should gather the observed properties of each binary from
their respective papers. You should then download and install
\href{https://cosmic-popsynth.github.io}{COSMIC} and check out the
documentation. In particular, you should find the section that shows how to
restart a binary and gives a specific example for Gaia BH1. Once you are able to
produce a population of black hole binaries, you shoud explore the degeneracies
between the pre-explosion separation, the strength and direction of the black
hole natal kick, and the final orbital period and eccentricity. Finally, you
should repeat the excersize with Gaia BH2's parameters. Do Gaia BH1 and Gaia BH2
have similar kick strengths?


\section{The Fastest Stars in the Galaxy}
Type Ia supernovae are extremely important in both astronomy and cosmology
because they are \emph{standardizable} candles. This means that with some
calibration, their brightness can be used to infer their brightness. This
technique was famously applied to measure the expansion of the Universe which
was found to be accelerating rather than just expanding. This finding led to
Saul Perlmutter, Brian Schmidt, and Adam Riess being awarded the Nobel Prize in
2011. 

Today, one of the main concerns about Type Ia supernovae is whether they are
caused from the mergers of two white dwarfs (double degenerate) or due to
accretion onto a white dwarf (single degenerate) or some combination of both!
One of the favored channels is the D6 scenario whose name is based on the
dynamically-driven double detonation of double degenerate (6 d's!) binaries. For
this project, you should first do a literature search to understand what the D6
scenario is and what observable signatures a D6 source would lead to.
\href{https://ui.adsabs.harvard.edu/abs/2018ApJ...865...15S/abstract}{Shen+2019}
is a good starting place, but there are lots of good references in Section 2 of
the paper. Once you have a handle on D6, use Gaia DR3 to find D6 candidates
based on the observable signatures of the channel. By the end of this project,
you should have a list of candidates and be able to argue why you think they are
or are not originating from the D6 channel. 

\section{Gravitational Wave Data Analysis}
The discovery of Gravitational Waves by the LIGO-Virgo Collaboration in 2015
marked a feat of science nearly 100 years in the making following their
prediction by Einstein's theory of General Relativity. The detection of two
merging black holes that were more than a billion light years away led to the
2016 Nobel Prize being awarded to Rai Weiss, Kip Thorne, and Barry Barish. The
work was carried out by \emph{over a thousand} people though -- the detection
was a physics experiemnt in and of itself!

Luckily for us, the LIGO Science Collaboration has produced tutorials for
interested folks to download and play with their data so that you can learn
exactly how gravitational wave data analysis works. For this project, you should
work through the "Guide to GW detections and noise" tutorial on the
\href{https://gwosc.org/tutorials/}{GW open science center website.}  Once
you've finished that tutorial, you can move onto the "Searching for
astrophysical sources" section to find an insprialing source. By the end of the
project, you should have a lot of figures and be able to describe how each
analysis step works.

\section{How old are those stars ... dynamically speaking?}
It is often \emph{very} hard to tell whether a star is very old based on
photometry alone since main sequence lifetimes for your average star are several
billion years. One way that we can measure the age is through kinematics because
old, low-metallicity stars are empirically known to orbit the Galaxy at larger
vertical distances from the center of the disk than younger, higher-metallicity
stars. In this project, you will explore how well kinematic ages approximate the
true age and metallicity of stars.  

First, you should download
\href{https://gala.adrian.pw/en/latest/index.html}{Gala} and work through the
tutorial which shows how to ingegrate orbits of stars observed by Gaia. Once you
have worked through the tutorial, you should gather a set of stars with known
metallicities and integrate their orbits. Do the orbits look like you'd expect
based on the metallicity? If not, what are some reasons that there could be
discrepancies. Based on the current position of the stars, is it safe to only
use the position as an age indicator? How much information does the addition of
the proper motion add?


\section{Exoplanet Orbital Dynamics with Rebound}
The discovery of over 5,000 exoplanets has made it clear that our own solar
system is not the only configuration for which planets can orbit stars on
long-term stable orbits. One of the biggest surprises of the 2010's was that
Jupiter-mass planets were found to orbit their host stars closer than Mercury
orbits in our own solar system. These planets were named "Hot Jupiters" due to
their proximity to their host star and were subsequently the subject of many
theoretical papers which tried to determine whether these planets could have
been scattered through gravitational interactions with other planets, slowly
migrated inward over time, or just been born there in the first place. 

For this project you should first do a literature search in ADS to trace the
reasoning for different theories and the evidence that backs each of them. You
should then download \href{https://rebound.readthedocs.io/en/latest/}{REBOUND},
an open-source code that you can use with Python or C and calculates the
evolution of gravitationally bound systems extremely acurately. There are
several examples for how REBOUND works for planet-star systems which you should
work through. After this, you should modify the examples to try to understand
how different assumptions for Hot Jupiter formation play out when integrating
different initial configurations of exoplanet systems. 

\end{document}